% Part: lambda-calculus
% Chapter: church-rosser
% Section: definitions-and-properties

\documentclass[../../../include/open-logic-section]{subfiles}

\begin{document}

\olfileid{lam}{cr}{dap}

\olsection{تعریف و خواص}

در این فصل مفهوم خاصیت چرچ--راسر و برخی خواص رایج مرتبط با آن را
معرفی می‌کنیم.

\begin{defn}[خاصیت چرچ--راسر، CR] رابطهٔ~$\xredone$ روی ترم‌ها را
  دارای \emph{خاصیت چرچ--راسر} می‌نامیم اگر و تنها اگر هرگاه
  $M \xredone P$ و~$M \xredone Q$، ترمی چون~$N$ وجود داشته باشد که
  $P \xredone N$ و~$Q \xredone N$.
\end{defn}

می‌توان حساب لامبدا را مدلی از محاسبه دانست که در آن ترم‌های دارای
صورت نرمال «مقدار»ها و یک رابطهٔ کاهش‌پذیری روی ترم‌ها «قواعد محاسبه»
هستند. خاصیت چرچ--راسر می‌گوید حتی هنگامی که برای ادامهٔ یک محاسبه بیش
از یک راه وجود دارد، عبارت همچنان فقط یک مقدار دارد.

برای نمونه‌ای از جبر مقدماتی، بیش از یک راه برای محاسبهٔ
$4 \times (1+2) + 3$ وجود دارد. می‌توان آن را به~$4 \times 3+3$ کاهش
داد (اگر نخست~$1+2$ را به~$3$ کاهش دهیم)، یا با استفاده از خاصیت
توزیع‌پذیری، نخست~$4 \times (1+2)$ را به~$4 \times 1+4 \times 2$
تبدیل کرد و به~$4 \times 1+4 \times 2+3$ رسید. بااین‌حال، هر دو عبارت
را می‌توان بیش از این به~$12+3$ کاهش داد.

اگر~$\xredone$ را کاهش~$\beta$ بگیریم، به‌آسانی می‌بینیم که یکی از
پیامدهای خاصیت چرچ--راسر این است که اگر ترمی صورت نرمال داشته باشد، آن
صورت یکتاست. فرض کنید~$M$ به~$P$ و~$Q$ کاهش یابد و هر دو صورت نرمال
باشند. بنا بر خاصیت چرچ--راسر، ترمی چون~$N$ وجود دارد که هم~$P$ و
هم~$Q$ به آن کاهش می‌یابند. چون بنا بر فرض~$P$ و~$Q$ صورت نرمال‌اند،
کاهش آن‌ها به~$N$ تنها می‌تواند کاهش بدیهی باشد؛ یعنی~$P$، $Q$ و~$N$
یکسان‌اند. این نتیجه سخن‌گفتن از \emph{صورت نرمال} یک ترم را موجه
می‌کند.

پس اگر حساب لامبدا را مدلی از محاسبه بدانیم، صورت نرمال یک ترم را
می‌توان «نتیجهٔ نهایی» محاسبه‌ای دانست که از آن ترم آغاز می‌شود. نتیجهٔ
بالا می‌گوید چنین محاسبه‌ای، اگر اصلاً نتیجهٔ نهایی داشته باشد، فقط یک
نتیجهٔ نهایی دارد؛ درست همان‌گونه که محاسبهٔ~$4 \times (1+2)+3$ فقط یک
نتیجه، یعنی~$15$، دارد.

\begin{thm} \ollabel{thm:str}
  اگر رابطهٔ~$\xredone$ خاصیت چرچ--راسر را داشته باشد و~$\xred$ کوچک‌ترین
  رابطهٔ تعدیِ شامل~$\xredone$ باشد، آنگاه~$\xred$ نیز خاصیت چرچ--راسر
  را دارد.
\end{thm}

\begin{proof}
  فرض کنید~$M \xred P$ و~$M \xred Q$. پس زنجیره‌هایی به‌صورت زیر وجود
  دارند:
  \begin{align*}
    M & \xredone P_1 \xredone \dots \xredone P_m \text{ و}\\
    M & \xredone Q_1 \xredone \dots \xredone Q_n.
  \end{align*}
  که در آن~$P_m=P$ و~$Q_n=Q$. قضیه را با ساختن شبکه‌ای~$N$ از ترم‌ها
  با ارتفاع~$m + 1$ و عرض~$n + 1$ ثابت می‌کنیم. با~$N_{i,j}$ ترمِ سطر
  $i$ام و ستون~$j$ام را نشان می‌دهیم.
  
  شبکهٔ~$N$ را چنان می‌سازیم که~$N_{i,j} \xredone N_{i+1,j}$ و
  $N_{i,j} \xredone N_{i,j+1}$. تعریف آن چنین است:
  \begin{align*}
    N_{0,0} &= M \\
    N_{i,0} &= P_i && \text{اگر } 1 \le i \le m \\
    N_{0,j} &= Q_j && \text{اگر } 1 \le j \le n \\
  \intertext{و در دیگر موارد:}
    N_{i,j} &= R
  \end{align*}
  که در آن~$R$ ترمی است که~$N_{i-1,j} \xredone R$ و~$N_{i,j-1}
  \xredone R$. بنا بر خاصیت چرچ--راسرِ~$\xredone$، چنین ترمی همواره
  وجود دارد.

  اکنون~$N_{m,0} \xredone \dots \xredone N_{m,n}$ و~$N_{0,n}
  \xredone \dots \xredone N_{m,n}$ را داریم. توجه کنید که~$N_{m,0}$
  همان~$P$ و~$N_{0,n}$ همان~$Q$ است. نتیجه از تعریف~$\xred$ به دست
  می‌آید.
\end{proof}

\end{document}
